%% 毕业设计小结

\begin{designsummary}
本次毕业设计围绕无人机具身时空推理中的双认知能力评测与优化展开。课题从早期的视频时序定位任务逐步扩展为包含图像空间推理、视频时间定位、自我状态认知和环境情况认知的综合测试基准，并进一步完成了面向证据一致性的训练集构建与模型优化实验。通过这一过程，我对数据集构建、仿真场景采集、多模态大模型评测、结构化输出解析、模型后训练和科研论文写作都有了更完整的理解。

在工程实现方面，本课题最大的收获是认识到高质量测试基准并不是简单汇总样本，而是需要从场景资产、任务定义、标注质量、评测协议和可视化展示等多个环节同时保证一致性。UAV-DualCog 的四阶段构建流程让我更加熟悉了复杂工具链设计，也让我意识到统一配置、统一路径、统一输出契约对可复现研究的重要性。

在研究认识方面，实验结果表明当前多模态大模型虽然具备一定语义理解能力，但在无人机场景中仍缺少稳定、均衡、可证据验证的双认知能力。尤其是自我状态认知弱于环境情况认知、语义判断与空间或时间证据脱节等现象，使我进一步理解了测试基准对于发现模型能力边界的价值。能力优化实验进一步说明，图像侧证据感知训练能够显著提升答案与目标框的一致性，并能迁移到部分视频自我行为定位任务，但环境可见性推理仍需要更强的时序整合能力。这也使我更清楚地认识到无人机多模态系统需要同时发展自我状态认知、环境情况认知和跨媒介证据协同能力。
\end{designsummary}
